\subsection{Factorization}\label{sec:mathematical_underpinnings:factorization}

Factorization is the process of breaking down a \textbf{number or expression} into a product of its factors. These factors, when multiplied together, will result in the original number or expression. The concept of factorization applies various mathematical objects, including integers and polynomials.\footnote{A polynomial is a mathematical expression with variables and coefficients, that involves only the operations of addition, subtraction, multiplication and exponentiation to nonnegative integer powers, and has a finite number of terms. For instance, $(x^2 - 4x + 7)$ is a polynomial of degree $2$, $(x - 2)$ is a polynomial of degree $1$, and $(1)$ is a constant polynomial of degree $0$ --- more information on polynomials can be found at \url{https://mathworld.wolfram.com/Polynomial.html}.}

\paragraph*{Integer Factorization.} Factoring an integer means expressing it as a product (i.e., the result of the multiplication) of two or more smaller positive integers greater than 1 (usually, prime numbers). For instance, the factorization of $12$ using only prime numbers is $2 \times 2 \times 3$. An integer which cannot be further factorized is called \textbf{prime} number.

\paragraph*{Polynomial Factorization.} Factoring a polynomial means expressing it as a product (i.e., the result of the multiplication) of two or more smaller --- that is, lower-degree --- non-constant\footnote{In the context of polynomial factorization, we do not consider constant polynomials such as $(1)$, $(2)$, or $(3)$ --- they are like the number $1$ in integer factorization, that is, irrelevant.} polynomials \textbf{with the same kind of coefficients} (e.g., integers, rationals, reals, or complex). For instance, the factorization of $(x^2 - 4)$ is $(x - 2)(x + 2)$. A polynomial which cannot be further factorized is called \textbf{irreducible}.

\remarkBlock{Integer vs. polynomial factorization}{Although factoring methods differ, integer and polynomial factorizations are logically equivalent: both suit the definition of writing a number (integer) or expression (polynomial) into a product of its factors. In simpler words, \textbf{for a polynomial to be irreducible is equivalent for an integer to be prime}.}

